% titel: Blatt 8, A1: Siedepunkt und intermolekulare Kräfte
% skript: Mortimer Kap. 12 (intermolekulare Kräfte)

% Teil I: Werkzeuge, kein Lösungsweg.

Der Siedepunkt einer Substanz hängt davon ab, wie stark die Teilchen im
flüssigen Zustand aneinander gebunden sind: Je stärker die intermolekularen
Kräfte, desto mehr Wärmeenergie braucht es, um die Teilchen in den
Gaszustand zu überführen.

\begin{werkzeugkasten}[Intermolekulare Kräfte und Siedepunkt]
\begin{itemize}
  \item \textbf{Dispersion {}= (Van-der-Waals-)Kräfte:} in allen Molekülen
        (auch in Edelgasen), wachsen mit der Elektronenzahl und damit mit der
        Molekülmasse.
  \item \textbf{Dipol-Dipol-Kräfte:} treten auf, wenn Moleküle ein permanentes
        Dipolmoment besitzen (polare kovalente Bindung).
  \item \textbf{Wasserstoffbrücken:} besonders starke Kräfte zwischen einem
        H-Atom und einem stark elektronegativen Atom (N, O, F).
  \item \textbf{Ionenkräfte:} bei Ionenverbindungen (z.~B. Salzen) wirken
        elektrostatische Coulomb-Kräfte.
\end{itemize}
\medskip
Reihenfolge (steigende zwischenmolekulare Anziehung $\rightarrow$ steigender
Siedepunkt): Dispersion < Dipol-Dipol < Wasserstoffbrücken < ionisch.
\textit{Skript:} Mortimer Kap. 12, Abschnitt zu intermolekularen Kräften.
\end{werkzeugkasten}

\begin{achtung}[Darauf achten]
\begin{itemize}
  \item Für \textbf{ionische} Stoffe (hier BaCl$_2$) entscheidet die Größe der
        Coulomb-Anziehung, nicht die Molmasse.
  \item Eine \textbf{polare} kovalente Bindung heisst noch nicht, dass ein
        Molekül ein Dipolmoment hat (Symmetrie! Bei 2. der Teilaufgaben wichtig).
  \item Kohlenmonoxid CO und Neon Ne sind beide kleine Moleküle, aber CO ist
        polar, Ne unpolar --- das beantwortet die Reihenfolge.
  \item Die Molmasse ist ein Hinweis, aber nicht allein entscheidend: HF hat
        trotz kleiner Molmasse einen hohen Siedepunkt wegen H-Brücken.
\end{itemize}
\end{achtung}
